\labsection{3}{Supervised workflow, evaluation, and model selection}{sec:lab03-supervised-workflow}

\begin{labproject}{From a no-skill baseline to a locked final model}
\textbf{Coverage.} Chapters 6--8.\\
\textbf{Core data.} Titanic and Breast Cancer Wisconsin. Heart Failure is an optional transfer activity inside the same lab.\\
\textbf{Goal.} Compare supervised baselines fairly, inspect threshold-dependent errors and ranking metrics, tune only on development data, and preserve an untouched final test boundary.

\begin{labmode}
This is the course's main workflow lab. Do not treat baseline comparison, threshold policy, cross-validation, and final evaluation as separate labs; they are successive decisions in one supervised-learning process.
\end{labmode}

\medskip\noindent\textbf{Part A --- Supervised baselines on Titanic.}\par
\begin{lstlisting}[style=mlpython]
from pathlib import Path
import pandas as pd
from sklearn.compose import ColumnTransformer
from sklearn.ensemble import RandomForestClassifier
from sklearn.impute import SimpleImputer
from sklearn.linear_model import LogisticRegression
from sklearn.metrics import accuracy_score, roc_auc_score
from sklearn.model_selection import train_test_split
from sklearn.pipeline import Pipeline
from sklearn.preprocessing import OneHotEncoder, StandardScaler
\end{lstlisting}

\begin{lstlisting}[style=mlpython]
PROJECT_ROOT = Path.cwd()
train_path = PROJECT_ROOT / "all_datasets" / "titanic_dataset" / "Titanic-Dataset.csv"
df = pd.read_csv(train_path)
\end{lstlisting}

\begin{lstlisting}[style=mlpython]
features = ["Pclass", "Sex", "Age", "SibSp", "Parch", "Fare", "Embarked"]
X = df[features]
y = df["Survived"]

num_cols = ["Age", "SibSp", "Parch", "Fare"]
cat_cols = ["Pclass", "Sex", "Embarked"]
preprocess = ColumnTransformer([
    ("num", Pipeline([
        ("imputer", SimpleImputer(strategy="median")),
        ("scale", StandardScaler()),
    ]), num_cols),
    ("cat", Pipeline([
        ("imputer", SimpleImputer(strategy="most_frequent")),
        ("onehot", OneHotEncoder(handle_unknown="ignore")),
    ]), cat_cols),
])
\end{lstlisting}

\begin{lstlisting}[style=mlpython]
X_train, X_valid, y_train, y_valid = train_test_split(
    X, y, test_size=0.25, random_state=42, stratify=y
)

models = {
    "logistic": LogisticRegression(max_iter=1000),
    "random_forest": RandomForestClassifier(
        n_estimators=300, random_state=42, min_samples_leaf=2
    ),
}
\end{lstlisting}

\begin{lstlisting}[style=mlpython]
for name, estimator in models.items():
    pipe = Pipeline([("preprocess", preprocess), ("model", estimator)])
    pipe.fit(X_train, y_train)
    pred = pipe.predict(X_valid)
    prob = pipe.predict_proba(X_valid)[:, 1]
    print(name, "accuracy=", round(accuracy_score(y_valid, pred), 3),
          "roc_auc=", round(roc_auc_score(y_valid, prob), 3))
\end{lstlisting}

\medskip\noindent\textbf{Part B --- Thresholds and ranking metrics.}\par
\begin{lstlisting}[style=mlpython]
from pathlib import Path
import pandas as pd
import matplotlib.pyplot as plt
from sklearn.linear_model import LogisticRegression
from sklearn.metrics import (
    ConfusionMatrixDisplay,
    average_precision_score,
    classification_report,
    confusion_matrix,
    roc_auc_score,
    RocCurveDisplay,
    PrecisionRecallDisplay,
)
from sklearn.model_selection import train_test_split
from sklearn.pipeline import Pipeline
from sklearn.preprocessing import StandardScaler
\end{lstlisting}

\begin{lstlisting}[style=mlpython]
PROJECT_ROOT = Path.cwd()
train_path = PROJECT_ROOT / "all_datasets" / "breast_cancer_dataset" / "data.csv"
df = pd.read_csv(train_path)
\end{lstlisting}

\begin{lstlisting}[style=mlpython]
X = df.drop(columns=["id", "diagnosis", "Unnamed: 32"], errors="ignore")
y = (df["diagnosis"] == "M").astype(int)

X_train, X_valid, y_train, y_valid = train_test_split(
    X, y, test_size=0.25, random_state=42, stratify=y
)
model = Pipeline([
    ("scale", StandardScaler()),
    ("model", LogisticRegression(max_iter=2000)),
])
model.fit(X_train, y_train)
prob = model.predict_proba(X_valid)[:, 1]
\end{lstlisting}

\begin{lstlisting}[style=mlpython]
# Explore threshold policies on development validation data, not a final test set.
for threshold in [0.50, 0.30]:
    pred = (prob >= threshold).astype(int)
    print(f"\nthreshold={threshold:.2f}")
    print(confusion_matrix(y_valid, pred))
    print(classification_report(y_valid, pred, digits=3))
    ConfusionMatrixDisplay.from_predictions(y_valid, pred)
    plt.title(f"Confusion matrix at threshold {threshold:.2f}")
    plt.tight_layout()
    plt.show()

print("ROC AUC:", round(roc_auc_score(y_valid, prob), 3))
print("PR AUC:", round(average_precision_score(y_valid, prob), 3))
\end{lstlisting}

\begin{lstlisting}[style=mlpython]
RocCurveDisplay.from_predictions(y_valid, prob)
plt.title("ROC curve")
plt.tight_layout()
plt.show()
\end{lstlisting}

\begin{lstlisting}[style=mlpython]
PrecisionRecallDisplay.from_predictions(y_valid, prob)
plt.title("Precision-recall curve")
plt.tight_layout()
plt.show()
\end{lstlisting}

\medskip\noindent\textbf{Part C --- Optional heart-failure transfer.}\par
\begin{lstlisting}[style=mlpython]
from pathlib import Path
import matplotlib.pyplot as plt
import pandas as pd
from sklearn.compose import ColumnTransformer
from sklearn.impute import SimpleImputer
from sklearn.linear_model import LogisticRegression
from sklearn.metrics import (
    ConfusionMatrixDisplay,
    PrecisionRecallDisplay,
    RocCurveDisplay,
    average_precision_score,
    classification_report,
    roc_auc_score,
)
from sklearn.model_selection import train_test_split
from sklearn.pipeline import Pipeline
from sklearn.preprocessing import OneHotEncoder, StandardScaler
\end{lstlisting}

\begin{lstlisting}[style=mlpython]
PROJECT_ROOT = Path.cwd()
train_path = PROJECT_ROOT / "all_datasets" / "heart_failure_dataset" / "heart.csv"
df = pd.read_csv(train_path)
\end{lstlisting}

\begin{lstlisting}[style=mlpython]
X = df.drop(columns="HeartDisease")
y = df["HeartDisease"]

num_cols = X.select_dtypes(include="number").columns.tolist()
cat_cols = X.select_dtypes(exclude="number").columns.tolist()
preprocess = ColumnTransformer([
    ("num", Pipeline([
        ("imputer", SimpleImputer(strategy="median")),
        ("scale", StandardScaler()),
    ]), num_cols),
    ("cat", Pipeline([
        ("imputer", SimpleImputer(strategy="most_frequent")),
        ("onehot", OneHotEncoder(handle_unknown="ignore")),
    ]), cat_cols),
])
\end{lstlisting}

\begin{lstlisting}[style=mlpython]
X_train, X_valid, y_train, y_valid = train_test_split(
    X, y, test_size=0.25, random_state=42, stratify=y
)
model = Pipeline([
    ("preprocess", preprocess),
    ("model", LogisticRegression(max_iter=2000)),
])
model.fit(X_train, y_train)
prob = model.predict_proba(X_valid)[:, 1]
pred = (prob >= 0.50).astype(int)

print(classification_report(y_valid, pred, digits=3))
print("ROC AUC:", round(roc_auc_score(y_valid, prob), 3))
print("PR AUC:", round(average_precision_score(y_valid, prob), 3))
\end{lstlisting}

\begin{lstlisting}[style=mlpython]
ConfusionMatrixDisplay.from_predictions(y_valid, pred)
plt.title("Heart-disease confusion matrix")
plt.tight_layout()
plt.show()
\end{lstlisting}

\begin{lstlisting}[style=mlpython]
RocCurveDisplay.from_predictions(y_valid, prob)
plt.title("Heart-disease ROC curve")
plt.tight_layout()
plt.show()
\end{lstlisting}

\begin{lstlisting}[style=mlpython]
PrecisionRecallDisplay.from_predictions(y_valid, prob)
plt.title("Heart-disease precision-recall curve")
plt.tight_layout()
plt.show()
\end{lstlisting}

\medskip\noindent\textbf{Part D --- Cross-validated model selection.}\par
\begin{lstlisting}[style=mlpython]
from pathlib import Path
import pandas as pd
from sklearn.compose import ColumnTransformer
from sklearn.impute import SimpleImputer
from sklearn.linear_model import LogisticRegression
from sklearn.metrics import roc_auc_score
from sklearn.model_selection import GridSearchCV, StratifiedKFold, train_test_split
from sklearn.pipeline import Pipeline
from sklearn.preprocessing import OneHotEncoder, StandardScaler
\end{lstlisting}

\begin{lstlisting}[style=mlpython]
PROJECT_ROOT = Path.cwd()
train_path = PROJECT_ROOT / "all_datasets" / "titanic_dataset" / "Titanic-Dataset.csv"
df = pd.read_csv(train_path)
\end{lstlisting}

\begin{lstlisting}[style=mlpython]
features = ["Pclass", "Sex", "Age", "SibSp", "Parch", "Fare", "Embarked"]
X = df[features]
y = df["Survived"]

num_cols = ["Age", "SibSp", "Parch", "Fare"]
cat_cols = ["Pclass", "Sex", "Embarked"]
preprocess = ColumnTransformer([
    ("num", Pipeline([
        ("imputer", SimpleImputer(strategy="median")),
        ("scale", StandardScaler()),
    ]), num_cols),
    ("cat", Pipeline([
        ("imputer", SimpleImputer(strategy="most_frequent")),
        ("onehot", OneHotEncoder(handle_unknown="ignore")),
    ]), cat_cols),
])
\end{lstlisting}

\begin{lstlisting}[style=mlpython]
pipe = Pipeline([
    ("preprocess", preprocess),
    ("model", LogisticRegression(max_iter=2000)),
])

X_train, X_test, y_train, y_test = train_test_split(
    X, y, test_size=0.25, random_state=42, stratify=y
)

cv = StratifiedKFold(n_splits=5, shuffle=True, random_state=42)
search = GridSearchCV(
    pipe,
    param_grid={
        "model__C": [0.01, 0.1, 1.0, 10.0, 100.0],
        "model__class_weight": [None, "balanced"],
    },
    scoring="roc_auc",
    cv=cv,
    n_jobs=-1,
)
search.fit(X_train, y_train)
\end{lstlisting}

\begin{lstlisting}[style=mlpython]
# Only after selection is complete do we evaluate once on the untouched test set.
test_prob = search.predict_proba(X_test)[:, 1]
print("best parameters:", search.best_params_)
print("training-CV ROC AUC:", round(search.best_score_, 3))
print("final test ROC AUC:", round(roc_auc_score(y_test, test_prob), 3))
\end{lstlisting}

\begin{decisionmemo}
Choose a classification threshold for a setting in which false negatives cost substantially more than false positives. State the consequence of lowering the threshold and name the validation evidence you used.
\end{decisionmemo}
\end{labproject}

\begin{labdeliverable}
Submit a comparison against a no-skill baseline, confusion-matrix and ROC/precision--recall evidence for at least two thresholds, the cross-validated selection result, and one final untouched-test result after all choices are locked. The Heart Failure transfer activity is optional and does not count as another lab.
\end{labdeliverable}
